% !TEX program = xelatex
\documentclass[sigconf,anonymous,nonacm]{acmart}
\usepackage{amsmath,amssymb}
\usepackage{graphicx}
\usepackage{booktabs,tabularx}
\graphicspath{{figures/}}
\settopmatter{printacmref=false,printfolios=true}
\setcopyright{none}
\acmConference[SAC '27]{The 42nd ACM/SIGAPP Symposium on Applied Computing}{April 5--9, 2027}{Gwangju, South Korea}
\acmDOI{}
\acmISBN{}
\hypersetup{hidelinks}
\setlength{\emergencystretch}{1.5em}
\newcommand{\artifacturl}{\url{https://anonymous.4open.science/r/nextphase-elimination-code-D1D9/}}
\raggedbottom
\begin{document}
\title[Next-Phase Elimination Risk in Battle Royale]{Predicting Next-Phase Elimination in Battle Royale: Target Alignment, Habitat Context, and Model Structure}
\begin{abstract}
We study next-phase team elimination from spectator-view PUBG telemetry through a controlled comparison of established prediction methods. From 13,948 matches, we construct 817,766 team--phase samples with a common feature cutoff; training, validation, and testing use successive game patches, and primary synthesis covers 11 cohorts. Reference gradient-boosted trees (GBDT) and a phase--team attention Transformer (PTR-Former) attain mean AUCs of 0.768 and 0.769, versus 0.747 for the phase-adjusted KRAFTON adaptation and 0.675 for the Sagong adaptation; PTR-Former's gap over KRAFTON persists among human teams (+0.020). Reconstructed KRAFTON inputs outperform ours within matches, and combining both input sets is best (GBDT AUC 0.786). PTR-Former improves within-phase but not within-match discrimination over GBDT. In the two largest cohorts, equal-trial tuning yields no consistent architecture winner, and truncating Cox learning at the next phase leaves no clear within-match difference from cross-entropy; cross-entropy's remaining within-phase advantage is consistent with match-level risk offsets, for which objectives normalized within a match's risk set provide no training signal. The survival adaptation still predicts winners better. Competition and remaining members retain risk associations after AI adjustment, whereas a movement residual is population-sensitive. These findings support aligning the prediction horizon, inputs, and evaluation unit before attributing gains to architecture.
\end{abstract}
\begin{CCSXML}
<ccs2012>
<concept><concept_id>10010147.10010257</concept_id><concept_desc>Computing methodologies~Machine learning</concept_desc><concept_significance>500</concept_significance></concept>
<concept><concept_id>10002951.10003227.10003351</concept_id><concept_desc>Information systems~Data mining</concept_desc><concept_significance>300</concept_significance></concept>
</ccs2012>
\end{CCSXML}
\ccsdesc[500]{Computing methodologies~Machine learning}
\ccsdesc[300]{Information systems~Data mining}
\keywords{spatio-temporal telemetry, discrete-time hazard, data leakage, survival analysis, esports analytics}
\maketitle
\section{Introduction}
\label{sec:introduction}
In a battle-royale game, teams must survive their opponents while the playable area shrinks. In PUBG: BATTLEGROUNDS (PUBG), up to 100 players compete individually or in teams of two or four. The safe zone shrinks in scheduled phases, and players outside it take continuous damage. At each phase boundary, a new safe zone is revealed. We study \emph{next-phase elimination risk}: the conditional probability that a team alive at that boundary is eliminated before the next phase ends.

The intended use is retrospective match review: an analyst compares teams' next-phase risks at a shared boundary from spectator-view telemetry, which joins all teams' trajectories with event and environment records~\cite{PUBGAPI,Zheng15}. Our inputs exclude the newly revealed zone and differ from a team's private information.

Risk depends on shrinking geometry and opponents, while its level changes across phases; correct rankings need not imply accurate probabilities. Later telemetry can reveal outcomes, so features, observed alive status, and relations must share an information cutoff.

Sagong et al.\ classify individual survival from activity scores~\cite{Sagong25a} or predict wins from strategy sequences~\cite{Sagong25b}. KRAFTON~\cite{KRAFTON} learns elimination-order risk with cross-team attention and converts it to win probability. These targets differ from next-phase elimination probability. Inputs, prediction units, and metrics also differ, so a performance gap alone does not identify the design choice responsible.

We compare these designs on a common landmark discrete-time task. A habitat-loss view~\cite{Fahrig03} organizes boundary exposure, movement~\cite{Nathan08}, competition, and members; it guides covariates rather than asserting an ecological mechanism.

\paragraph{Research questions.}
\begin{itemize}
\item \textbf{RQ1:} How are the four habitat factors associated with next-phase elimination risk, and are the associations consistent across phases and cohorts defined by map, perspective, and team size?
\item \textbf{RQ2:} How does a design learning next-phase elimination probability from habitat context compare with prior designs adapted to the same conditions?
\item \textbf{RQ3:} How do inputs, objective, phase baseline, and model structure change performance under the evaluated controls, and does this vary by cohort?
\end{itemize}

\paragraph{Contributions.}
Our central contribution is a matched-control measurement of a known property in landmark prediction: aligning the learning horizon shrinks the within-match AUC gap between Cox and cross-entropy below 0.001, whereas a within-phase gap of 0.006--0.008 remains and falls to 0.002 or less after centering within match--phases, consistent with a match-level risk offset that risk-set normalization gives no training signal (Sections~\ref{sec:results:rq3}, \ref{sec:discussion:drivers}); most of our designs' within-phase gain over phase-adjusted KRAFTON behaves the same way. Supporting it, we define common-cutoff, patch-separated samples with reconstruction checks (Section~\ref{sec:method}), compare established designs adapted to this task with a regularized logistic baseline (Section~\ref{sec:results:rq2}), and identify complementary inputs, metric-dependent structural effects, and habitat associations whose phase dependence and AI sensitivity matter more than their direction (Section~\ref{sec:results:rq1}). This is an empirical study, not a new hazard estimator.

\section{Related Work}
\label{sec:related}

\subsection{In-Game Outcome and Risk Prediction}
\label{sec:related:esports}
PUBG outcome prediction began with final placement from post-match statistics~\cite{KagglePUBG,LeeLee25}; other esports predict wins or short-horizon deaths from in-game information~\cite{Hodge21,Katona19,Marshall22}.
We adapt the PUBG activity-score~\cite{Sagong25a} and survival-learning~\cite{KRAFTON} designs to shared samples, preserving their intended inputs and training where telemetry permits (Table~\ref{tab:compare}).
\subsection{Survival Analysis, Ecology, and Multi-Agent Models}
\label{sec:related:survival}
The Cox model estimates relative risks without specifying a baseline hazard, the risk shared by all at a given time~\cite{Cox72}, whereas discrete-time survival models estimate interval event probabilities with interval baselines~\cite{SingerWillett93}, also with neural networks~\cite{Gensheimer19}; landmarking estimates subsequent risk for those alive at a prediction time~\cite{vanHouwelingen07}.
DeepHit combines event-time likelihood and ranking~\cite{DeepHit18}, and Dynamic-DeepHit learns from longitudinal measurements~\cite{DynamicDeepHit20}; our auxiliary placement-ranking term is not their event-time loss.
Because a baseline turns relative-risk scores into interval probabilities~\cite{PrenticeGloeckler78}, we measure the objective's effect that remains after attaching one.

Ecology studies how the amount and arrangement of habitat relate to population persistence~\cite{Fahrig03}; we use it only to construct and interpret risk covariates.
In multi-agent scenes, attention is combined across agent and time axes~\cite{Ngiam22} or graph structure is injected as attention bias~\cite{Ying21}, and player-relation graphs have been used in sports~\cite{Xenopoulos21}.
We measure the conditional contribution of such structure with inputs and objective fixed.

%======================================================================
\section{Data and Prediction Task}
\label{sec:method}

We construct one sample per team alive at a phase boundary and predict elimination before the next phase ends. Figure~\ref{fig:pipeline} connects data collection and sample construction to the three analyses and their evaluation.

\begin{figure*}[t]
\centering
\includegraphics[width=\textwidth]{fig_pipeline_clean.pdf}
\caption{Data collection, prediction-time sample construction, model comparisons, and evaluation. Panels A--B define the data and task; panels C--D organize the three research questions and their evaluation.}
\label{fig:pipeline}
\Description{Telemetry is converted to samples using a common cutoff, then analyzed for habitat associations, aligned model comparisons, and design decomposition.}
\end{figure*}

\subsection{Data Collection}
\label{sec:method:data}
PUBG's public telemetry records every player's position about every 10\,s, together with timestamped events such as damage, knock-downs, kills, items, safe-zone states, and phase changes~\cite{PUBGAPI}.
We queried the public API's list of sample matches from the past 24\,h 35 times between March 27 and June 2, 2026 (UTC), collecting Steam matches created from March 25 to June 1.
The collection spans PC patches 40.2--41.2~\cite{PUBGPatchNotes2026}, before the zone revamp~\cite{PUBGBlueZone2026}. Update 41.1 changed bot behavior and introduced Jammer Pack tiers with different Blue Zone protection, so a shared zone schedule does not imply a stationary prediction environment.
Prior processed data are unavailable~\cite{Sagong25a,KRAFTON}; we release our training, evaluation, and figure code with feature definitions (Section~\ref{sec:discussion:limits}).
We used only public match records, identified no individuals, and list data sizes in Table~\ref{tab:cohorts}.

\subsection{Cohorts and Match-Level Split}
\label{sec:method:cohort}
Matches with the same map, perspective (FPP/TPP), and team size form a cohort, since these change travel distances and fighting.
We keep normal and ranked matches, excluding tutorial and training matches, casual matches populated mostly by AI accounts, and the Vikendi map with many non-player combat accounts; because the game still adds AI accounts to normal matches, each team's AI-account share is an input (teams are never mixed, so it acts as an AI-team indicator).
We retain the 12 available cohorts and split complete matches chronologically by patch: 40.2 for training (4,439 matches), 41.1 for validation (6,971), and 41.2 for testing (2,538); validation is largest because our sampled matches span 35 days of patch 41.1, versus 13 for 40.2 and 16 for 41.2. Every model uses the same match manifest. Scheduled phase lengths and all preprocessing statistics are estimated on patch 40.2 alone, and features are rebuilt from raw telemetry under them~\cite{Kaufman12}. Before training, we required at least 50 training matches for primary cohort synthesis: eleven cohorts qualify (13,615 matches), and Karakin, with six training matches, is reported separately. Selection uses patch 41.1 and evaluation uses 41.2.

Of 145,019 primary test rows, 19,728 (13.6\%) are AI teams; elimination rates are 37.9\% for AI and 25.1\% for human teams, motivating human-only evaluation.

\begin{table}[t]
\centering
\caption{Matches, Team--Phase Samples, and Next-Phase Elimination Rate by Cohort}
\label{tab:cohorts}
\footnotesize
\setlength{\tabcolsep}{2.7pt}
\begin{tabular}{lccrrrr}
\hline
Map & View & Size & Matches & Train/val/test & Samples & Elim.\ (\%) \\
\hline
Erangel & TPP & 1 & 224 & 86/100/38 & 31,740 & 35.1 \\
Erangel & TPP & 2 & 656 & 238/308/110 & 59,170 & 29.5 \\
Taego & TPP & 2 & 573 & 214/246/113 & 47,223 & 29.2 \\
Erangel & FPP & 4 & 430 & 118/209/103 & 21,849 & 27.8 \\
Erangel & TPP & 4 & 3,283 & 1,094/1,627/562 & 196,442 & 25.7 \\
Karakin & TPP & 4 & 333 & 6/212/115 & 12,669 & 32.1 \\
Miramar & TPP & 4 & 1,543 & 521/733/289 & 92,279 & 24.6 \\
Paramo & TPP & 4 & 333 & 198/72/63 & 9,963 & 30.5 \\
Rondo & TPP & 4 & 1,039 & 306/613/120 & 53,683 & 25.2 \\
Sanhok & TPP & 4 & 940 & 222/564/154 & 48,793 & 27.6 \\
Taego & FPP & 4 & 435 & 118/226/91 & 23,390 & 30.0 \\
Taego & TPP & 4 & 4,159 & 1,318/2,061/780 & 220,565 & 24.9 \\
\hline
Total & & & 13,948 & 4,439/6,971/2,538 & 817,766 & 26.6 \\
\hline
\end{tabular}
\\[2pt]
\parbox{\columnwidth}{\scriptsize About 322.4\,M records (96.43\,M positions, 225.97\,M events); samples 2.9\,GB.}
\end{table}

\subsection{Time Axis and Information Cutoff}
\label{sec:method:samples}
Times are seconds since match start; phase $k$ of match $m$ runs from the reveal of the $k$-th safe zone to the end of the shrink to that zone.
We set the prediction time $b_{m,k}$ to the moment the shrink ends and the next safe zone is revealed, and build inputs only from records before the feature cutoff $a_{m,k}=b_{m,k}-\Delta$ with $\Delta=10$\,s (Fig.~\ref{fig:pipeline}B): if the fourth zone is revealed at 1,200\,s, the cutoff is 1,190\,s and a fight at 1,195\,s is excluded. Inputs, relations, and attention keys use records strictly before $a_{m,k}$; the 10\,s position grid takes each player's nearest record within $\pm5$\,s and admits only grid times $t$ with $t+5<a_{m,k}$.
A killed player counts as alive again only if a healthy position record appears at least 20\,s after the kill and before $a_{m,k}$.
We check the cutoff by rebuilding the inputs of sampled boundaries after deleting all records after $a_{m,k}$ and comparing them with the originals; a builder without the cutoff tests that the check detects leakage.

\subsection{Event-Based Labels and Estimand}
Knocked players can be revived and are not counted as eliminated; the elimination time $T_{m,i}$ of team $i$ in match $m$ is the kill time of its last member who did not return ($T_{m,i}=\infty$ for the winner, observed until the match end $C_m$).
The retrospective risk set is $\mathcal{R}_{m,k}=\{i:T_{m,i}>b_{m,k}\}$. Eligibility at $b_{m,k}$ is distinct from inputs fixed at $a_{m,k}$: teams finally eliminated in the intervening 10\,s are excluded, and retained rows are checked against alive status observed before the cutoff.
Each team's label is $Y_{m,i,k}=\mathbb{I}[T_{m,i}\le b_{m,k+1}]$, where $\mathbb{I}[\cdot]$ is 1 if its argument is true and 0 otherwise, and $b_{m,k+1}=C_m$ if phase $k+1$ is the last.
Samples $(m,i,k)$ are built at boundaries $k=1,\dots,7$ (in solo matches a team is a player), and the estimand is the next-phase elimination risk
\begin{equation}
\lambda_{m,i,k}=\Pr\!\left(Y_{m,i,k}=1 \mid i\in\mathcal{R}_{m,k},\,\mathcal{H}_{m}(a_{m,k})\right),
\label{eq:hazard}
\end{equation}
where $\mathcal{H}_{m}(t)$ is all telemetry of match $m$ recorded before time $t$.
This is an interval event probability~\cite{vanHouwelingen07}, not a per-second rate; events at $b_{m,k}$ leave the current risk set, and events at the right endpoint enter its label. Boundaries are excluded if match end is missing or a gap exceeding 60\,s in the pooled position stream overlaps the outcome window.

\section{Habitat Representation and Model Designs}
\label{sec:method:design}
\subsection{Habitat Context and Other Inputs}
\label{sec:method:representation}
As population persistence in a shrinking habitat depends on location and competition, we describe a team's state with four constructs of position and crowding.
\emph{Boundary exposure} is how far a team is outside the current zone (outside depth) or inside it (inside depth).
\emph{Movement residual} is the retrospective speed difference $(d-\ell)/\tau_k$~\cite{Nathan08}. Here $d$ is the nonnegative distance from the last observed team centroid to the boundary of the target zone already visible at the cutoff, $\ell$ sums horizontal displacements of the team centroid over alive grid points from phase start to the cutoff, and $\tau_k$ is its scheduled duration, estimated on training matches. It is an empirical proxy, not a future travel requirement: motion away from the zone also adds to $\ell$, and inside-zone teams ($d=0$) get a negative residual when they move. We also analyze its components, required speed $d/\tau_k$ and observed path speed $\ell/\tau_k$.
\emph{Local competition density} counts other teams within 200\,m (3-D centroid distance) at the focal team's last alive grid point, and \emph{remaining members} counts members alive at the cutoff.
Our inputs comprise a habitat block with the first three constructs (remaining members are in team status), other team-feature blocks (Table~\ref{tab:blocks}), and, for the Transformer of Sec.~\ref{sec:method:rq2}, the team centroid and three same-phase relations: proximity (dropped for pairs that fought), combat, and shared zone pressure (two teams within 200\,m, both outside or within 150\,m of the zone edge).

\begin{table}[t]
\centering
\caption{Input Blocks and Representative Features}
\label{tab:blocks}
\footnotesize
\setlength{\tabcolsep}{3pt}
\begin{tabular}{p{1.45cm}p{5.8cm}r}
\hline
Block & Representative features & \# \\
\hline
\multicolumn{3}{l}{\emph{Ours}} \\
\textbf{Habitat context} & \textbf{outside/inside depth, distance to boundary, required speed and movement residual, enemy teams within 100--500\,m, zone size and shrink rate, entry congestion} & 44 \\
Match context & phase index, phase length, team size, alive teams & 6 \\
Team status & alive members, mean/min health, member spread & 9 \\
Movement & path length, mean speed, tortuosity, zone-entry time & 7 \\
Combat & damage dealt/taken, kills/knocks, opponents engaged & 13 \\
Resources & item pickup/equip/use, healing, vehicle use & 16 \\
Spatial/elev. & nearest-enemy distance, enemy density, elevation & 12 \\
Relation sum. & degree and strength of proximity/combat/pressure & 43 \\
Relation graph & proximity, combat (out/in), shared-pressure weights & --- \\
AI indicator & AI-account share (0 or 1), appended to our tabular and neural inputs & 1 \\
\hline
\multicolumn{3}{l}{\emph{Prior designs}} \\
Sagong et al. & phase-wise combat, distance, density, and item counts; fight/activity scores; play type; type-wise survival & 12 \\
KRAFTON & kills/knocks, weapon damage, damage taken, healing, equipment, nearby allies/enemies, zone distances, health, alive members, positions & 55+ \\
\hline
\end{tabular}
\smallskip
\parbox{\columnwidth}{\scriptsize The first eight numeric blocks contain 150 features. Reference GBDT adds the AI indicator (151 inputs); matched tuning also includes the two centroid coordinates (153). Original-input Sagong and KRAFTON adaptations have no explicit AI indicator.}
\end{table}

\subsection{RQ1: Interpretable Hazard Model}
\label{sec:method:rq1}
For RQ1 we fit, per cohort, a directly interpretable logistic regression:
\begin{equation}
\mathrm{logit}\,\lambda_{m,i,k}=\alpha_k+\textstyle\sum_{c\in\mathcal{C}}\beta_c z_{c,m,i,k},
\label{eq:ehm}
\end{equation}
where $\mathrm{logit}\,p=\log\{p/(1-p)\}$ and $\alpha_k$ is the phase intercept, representing the risk level shared by its teams on the log-odds scale.
$\mathcal{C}$ contains five variables (outside depth, inside depth, movement residual, local competition density, and remaining members) and, as controls, the log number of alive teams and the AI-team indicator at the cutoff; solo cohorts drop remaining members.
$z_{c,m,i,k}$ is variable $c$ standardized by the cohort's training mean and SD, so $\beta_c$ is the change in log-odds per SD.
Standard errors are clustered by match, and the 11 eligible cohort coefficients are pooled with a random-effects model~\cite{DerSimonian86}, reporting 95\% prediction intervals for a new cohort~\cite{Higgins09}. Sensitivity analyses omit the AI indicator, restrict targets to human teams, vary the competition radius to 100/500 m, and replace the movement residual by its two components. Effects of one more enemy team or member are averaged probability changes from adding one to the physical covariate.
Phase bands (early 1--2, middle 3--5, late 6--7) get coefficients in within-band SD units, and their equality is tested with a multivariate random-effects Wald test, also on the common training-SD scale.
An exploratory lite model with phase intercepts and four variables (required speed, competition density, log alive teams, remaining members) measures how much of the discrimination ($\mathrm{AUC}-0.5$) of our full tree model (GBDT, Sec.~\ref{sec:method:rq2}) they reproduce; it drops both depth terms, as outside depth is nearly collinear with required speed, and keeps log alive teams as a control.

\subsection{RQ2: Prediction Designs}
\label{sec:method:rq2}
Our design learns the next-phase elimination probability directly, per phase:
\begin{equation}
\mathrm{logit}\,\lambda_{m,i,k}=\alpha_k+f(\mathbf{x}_{m,i,k},\mathcal{G}_{m,k}),
\label{eq:logit}
\end{equation}
where $\mathbf{x}_{m,i,k}$ is team $i$'s input feature vector including habitat context, $\mathcal{G}_{m,k}$ is the graph of the three relations among teams alive at $a_{m,k}$, and $f$ is a learned function.
Models minimize the cross-entropy $-\sum[Y_{m,i,k}\log\lambda_{m,i,k}+(1-Y_{m,i,k})\log(1-\lambda_{m,i,k})]$ over all training team--phase samples.
Equation~\eqref{eq:logit} is the conceptual form: our two learners represent $\alpha_k$ through phase inputs, and only one decomposition step adds explicit intercepts. The learners are gradient-boosted decision trees (GBDT)~\cite{Friedman01} that take the phase index as input, and PTR-Former (phase--team--relation Transformer), a reference model for decomposing structure.
PTR-Former alternates attention over the same team's earlier phases (phase axis) and over other teams alive in the same phase (team axis)~\cite{Vaswani17}, adds learned relation-type biases to team-axis attention~\cite{Ying21}, and adds a ranking loss on final placement (weight 0.1). As a same-target linear baseline, L2-regularized logistic regression uses phase indicators and the GBDT's 151 inputs standardized on training data; validation log loss selects the penalty from five values.

KRAFTON's predictor instead learns elimination ordering: its public weighted survival objective combines time-gap weights, eliminated-team log probabilities, and survivor-complement terms. It is not standard Cox partial likelihood, although it uses the risk-set term $\exp(s_i)/\sum_{j\in\mathcal{R}}\exp(s_j)$~\cite{Cox72}, where $s_i$ is team $i$'s score at a time point and $\mathcal{R}$ the set of teams alive then.
Both terms depend on scores only through these normalized ratios, and the time-gap weights depend on outcomes, so the full objective is invariant to a common shift within each risk set and identifies neither an absolute phase nor a match risk level. Before phase-specific mapping, the reference adaptation uses validation-fitted Platt scaling $p=\sigma(A s+B)$, where $\sigma(u)=(1+e^{-u})^{-1}$.
We therefore also compare the probability $1-\exp\{-\exp(\gamma_k+s_{m,i,k})\}$, which attaches a phase baseline to the score; $s_{m,i,k}$ is team $i$'s score at the cutoff $a_{m,k}$ and $\gamma_k$ is phase $k$'s baseline hazard~\cite{PrenticeGloeckler78}. The original-input adaptation fits $\gamma_k$ on validation risk sets; the our-input path fits it on training risk sets. All matched objective--mapping controls fit their mapping on validation data.

These are task adaptations (Table~\ref{tab:compare}). KRAFTON uses public commit 99ed9d1 with integration fixes; its 55 features include 32 definition-matched reconstructions, 22 proxies, and one zero-filled feature. Sagong is a paper-based reimplementation with attack/assist proxies, aggregated from its solo, half-phase task to teams and full phases. Neither reproduces original empirical performance; we also report KRAFTON's winner accuracy and elimination-order C-index.

\begin{table*}[t]
\centering
\caption{Prior Designs and Ours}
\label{tab:compare}
\footnotesize
\setlength{\tabcolsep}{3pt}
\begin{tabular}{p{0.10\textwidth}p{0.25\textwidth}p{0.30\textwidth}p{0.30\textwidth}}
\hline
 & Sagong et al.~\cite{Sagong25a} & KRAFTON~\cite{KRAFTON} & Ours \\
\hline
Unit & player & team & team (player in solo) \\
Target & next half-phase survival & elimination-order risk (win probability) & \textbf{next-phase elimination probability} \\
Prediction time & half-phase boundaries & many in-match time points & phase boundaries (next zone revealed) \\
Phase risk level & type-wise survival rates & not represented & \textbf{phase-conditioned} (phase as input) \\
Team interaction & none & cross-team attention & attention, relation bias (PTR-Former) \\
Objective & survival classification & weighted survival, Cox-like risk sets & \textbf{next-phase cross-entropy} \\
\hline
\end{tabular}
\end{table*}

\subsection{RQ3: Conditional Design Decomposition}
\label{sec:method:decomposition}
RQ3 asks where performance differences come from, so we change each prior design step by step toward ours and read each step's change as its contribution conditional on the preceding design.
On the Sagong path, with GBDT fixed, we add habitat context to Sagong's inputs (a1) and replace the rest with our inputs (a2); separately, (b) replaces KRAFTON's inputs with ours.
The KRAFTON path starts from the KRAFTON architecture trained with the upstream weighted survival objective on our inputs, then adds a phase baseline (c1), switches to next-phase cross-entropy (c2), and adds explicit phase intercepts (d); structure effects replace the architecture with a team-axis-only PTR-Former (e) and add the phase axis and relation bias (f), both without the ranking loss.
These are path-conditional comparisons: the first KRAFTON step changes inputs and training rows together, the objective switch also changes the head intercept initialization and stopping criterion, (e) changes architecture and training recipe, and (f) adds two components. The additional matched controls below address these confounds; path gains alone are not a factorial attribution.
Step (g) adds KRAFTON's inputs to GBDT to test complementarity.
Five primary comparisons were fixed before the patch split was defined: a 2$\times$2 on/off design of the phase and team axes (AUC for both axes and, for the team axis, Spearman rank correlation within matches), a control with relations shuffled within matches, and a zone token (the safe zone as an extra attention token) tested without habitat and zone features.

\section{Evaluation Protocol}
\label{sec:method:evaluation}
All designs are scored on identical test samples (Figure~\ref{fig:pipeline}D); RQ1 inference instead uses the clustered standard errors and Wald tests of Section~\ref{sec:method:rq1}.

\subsection{Discrimination and Probability Accuracy}
The primary metric, AUC, is the probability that an eliminated team receives a higher risk than a surviving one.
Because elimination rates vary by phase, predicting the phase alone can raise AUC, so we also report within-phase AUC (unweighted mean over phases with at least 30 rows and both outcome classes) and within-match AUC, which compares only teams alive in the same match at the same moment~\cite{HeagertyZheng05} (pairs pooled over risk sets containing both classes).
Same-match pairs are about $1/M$ of within-phase pairs for $M$ test matches, so within-phase AUC mostly ranks teams across matches. As a post hoc diagnostic added after the objective results were examined, we locate a within-phase difference by recomputing it after subtracting each model's mean logit per match--phase and by computing AUC over different-match pairs from these means alone; neither is an additive decomposition of AUC into causes. Probability accuracy uses log loss and Brier score. Phase calibration error is $\sum_k(n_k/n)|\bar p_k-\bar y_k|$, where $n_k$ is phase $k$'s sample count, $n$ the cohort total, and bars denote mean prediction and outcome. Within-phase ECE~\cite{Naeini15} further partitions each phase into ten equal-width probability bins: $\sum_{k,j}(n_{kj}/n)|\bar p_{kj}-\bar y_{kj}|$. Both errors are computed per cohort before averaging. For review use, top-3 capture is the share of next-phase eliminations among the three highest-risk teams at each boundary of a match (ties split at random; boundaries without eliminations excluded), compared with random selection and a fewest-remaining-members rule. Calibration slope and calibration-in-the-large are the coefficient of a logistic regression of outcomes on predicted log-odds and its intercept with that coefficient fixed at one~\cite{VanCalster16}.

\subsection{Statistical Comparisons}
Per-cohort confidence intervals come from a match-level paired bootstrap with 2,000 test-match resamples.
Cohorts are pooled with a DerSimonian--Laird random-effects model~\cite{DerSimonian86}, checked with REML and Hartung--Knapp intervals~\cite{HartungKnapp01}.
Decision rules, the SESOI, and the five primary comparisons were fixed before the patch-split models were trained. A difference is a gain if its one-sided $p<0.05$ and its estimate is at least the smallest effect size of interest (SESOI), 0.005 (0.01 for rank correlation; Section~\ref{sec:results:rq3} relates it to top-3 capture), practically equivalent if its 90\% CI lies within $\pm$SESOI~\cite{Schuirmann87}, and inconclusive otherwise.
Holm correction applies to the five one-sided superiority tests against zero~\cite{Holm79}; the 90\% equivalence intervals are unadjusted, and a gain is evidence against zero, not evidence that the true effect exceeds the SESOI (Table~\ref{tab:primary}). The patch protocol was recorded before training; equal-trial tuning (Table~\ref{tab:tuning}) selects settings on patch 41.1 alone and is our cleanest structural comparison.
\subsection{Training and Robustness Checks}
Reference GBDT uses 200 trees (learning rate 0.06); PTR (width 128, four heads, two blocks) and KRAFTON (width 64, eight heads/layers) use AdamW ($10^{-3}$/$10^{-4}$) with validation stopping, and the artifact lists all settings.

Equal-trial tuning~\cite{Bergstra12,Cawley10} compares GBDT, a KRAFTON-derived team-only backbone, and PTR without ranking loss in the two largest cohorts, each with the same 153 columns (150 features, centroid coordinates, AI indicator), BCE, and 16 configurations selected by validation NLL. Neural models share the search list, a 12-epoch cap, three-epoch patience, and eight-match batches; selected settings are rerun with seeds 42/13/2026, and metrics average the runs.

Additional controls hold observations, node inputs, and the backbone fixed to compare full-follow-up versus next-phase weighted survival objectives, standard Cox with Breslow ties at both horizons, and BCE. For weighted truncation, only events beyond the horizon are excluded; full-follow-up time weights are retained. Standard Cox uses administrative horizon censoring. All five conditions use a configuration fixed before tuning and validation within-phase AUC for stopping, followed by common validation logistic calibration with phase indicators. An explicit $2\times2$ analysis crosses objective (full-follow-up weighted survival/BCE) with probability mapping (global/phase-specific); the objective factor bundles loss and horizon. Its Shapley contributions average the two possible orders~\cite{Shorrocks13} and sum to the joint change in every bootstrap draw. Test data select neither settings nor mappings.

Human-team evaluations retain opponents as observed context; retraining controls remove the AI input or restrict the training loss to human teams, keeping the selected hyperparameters. History controls add summaries of earlier tokens to the current token, or retrain PTR on ordered versus shuffled past prefixes with the current token fixed.

An exploratory GATv2~\cite{Brody22} treats each match--phase as a graph with four continuous relation channels and uses BCE with the same 16 configurations and three seeds; a control jointly relabels relation channels. It was added after the reference results had been examined.


%======================================================================
\section{Results}
\label{sec:results}

Rebuilding inputs from telemetry truncated at the cutoff reproduced all 453 sampled cutoffs (72 matches, stratified by cohort and split), whereas disabling truncation changed 450 of them. Only two of the 817,766 retained rows, both in training, lacked an observed living member at the cutoff.
Unless stated otherwise, primary synthesis covers the 11 cohorts with at least 50 training matches; differences use random-effects pooling with 95\% CIs. Counts in parentheses give the number of evaluable cohorts with the pooled sign.
An AUC difference of 0.01 means a 1-percentage-point (pp) higher probability of ranking an eliminated team above a survivor; gains and equivalence follow Sec.~\ref{sec:method:evaluation}.




\subsection{RQ1: Habitat Context and Next-Phase Risk}
\label{sec:results:rq1}
\begin{figure}[t]
\centering
\includegraphics[width=\columnwidth]{fig_rq1_effects.pdf}
\caption{AI-adjusted habitat associations in training patch 40.2. Points and 95\% intervals pool 11 eligible cohorts (10 for members). Each band uses its own SD scale.}
\label{fig:rq1}
\Description{Forest plot of standardized habitat-factor coefficients and confidence intervals for early, middle, and late phases.}
\end{figure}
Competition and remaining members have the clearest all-phase associations in the AI-adjusted model. A one-SD increase in local competition corresponds to $+0.284$ $[0.230,\allowbreak 0.339]$ log-odds (11/11), and remaining members to $-0.594$ $[-0.686,\allowbreak -0.502]$ (10/10; solo excluded), or odds multipliers of approximately 1.33 and 0.55. Their signs agree across cohorts, but their magnitudes are heterogeneous ($I^2=0.90$ and 0.99); their prediction intervals, $[0.08,\allowbreak 0.49]$ and $[-0.95,\allowbreak -0.24]$, still exclude zero. Neither outside depth ($+0.002$ $[-0.021,\allowbreak 0.025]$) nor inside depth ($-0.020$ $[-0.054,\allowbreak 0.014]$) has a clear pooled all-phase association. The largest VIF is 2.47.

\paragraph{Movement interpretation depends on population and specification.}
The movement coefficient falls from $+0.069$ $[0.052,\allowbreak 0.086]$ without AI adjustment to $+0.005$ $[-0.007,\allowbreak 0.018]$ with adjustment and $+0.010$ $[-0.004,\allowbreak 0.023]$ in human teams, and it excludes zero only at a 500\,m competition radius; competition and member directions survive these checks. The residual is negative in a median 88.1\% of human and 93.8\% of AI rows ($d>0$ in 37.4\%/28.2\%), so it rarely describes an unmet travel requirement; separating required and observed speeds gives positive coefficients ($+0.395$/$+0.025$ per SD) but raises the maximum VIF to 12.55.

\paragraph{Associations vary across phases.}
In within-band SD units (Fig.~\ref{fig:rq1}), inside depth is near zero early, then $-0.224$/$-0.414$ middle/late (11/11). Competition is $+0.385$/$+0.213$/$+0.237$, and members $-0.516$/$-0.601$/$-0.617$. Equality tests reject constant depth and competition slopes on both SD scales ($p\le2.4\times10^{-6}$). Model-implied $+1$ enemy-team changes are 9.8/3.7/3.9\,pp; $+1$ member changes are $-6.6$/$-10.5$/$-13.3$\,pp.

\paragraph{A small model retains much of the discrimination.}
The lite model's mean overall/within-phase AUCs are 0.719/0.677, versus 0.615/0.500 for phase-only and 0.768/0.742 for full GBDT. Of the GBDT's AUC above chance (0.5), phase intercepts alone reproduce 43.1\% $[41.4,\allowbreak 44.8]$ and the four-variable lite model 81.9\% $[80.5,\allowbreak 83.1]$. By within-phase AUC, the lite model retains 73.3\% $[71.2,\allowbreak 75.7]$. Restoring both depth terms adds 0.0034 $[0.0023,\allowbreak 0.0044]$ overall and 0.0068 $[0.0054,\allowbreak 0.0080]$ within phases. These descriptive models omit the AI indicator.

\subsection{RQ2: Comparison with Prior Designs}
\label{sec:results:rq2}
\begin{figure*}[t]
\centering
\includegraphics[width=\textwidth]{fig_contrast.pdf}
\caption{(a) Observed and predicted next-phase elimination rates (11-cohort means). (b) Matched objective controls in the two largest cohorts. (c) Within-phase AUC differences before and after match--phase centering, with within-match differences; KRAFTON+phase is phase-adjusted KRAFTON. Horizontal bars indicate 95\% confidence intervals; pooled rows use random-effects estimates across 11 cohorts.}
\label{fig:contrast}
\Description{Per-phase observed and predicted elimination rates; AUC differences of cross-entropy from full-follow-up and next-phase Cox within matches and within phases; within-phase differences before and after match-phase centering for several design contrasts.}
\end{figure*}
\begin{table*}[t]
\centering\footnotesize
\caption{Patch-41.2 performance: arithmetic means over 11 eligible cohorts. Human AUC restricts evaluation to human teams; other columns use all teams.}
\label{tab:main}
\setlength{\tabcolsep}{2.5pt}
\begin{tabular*}{\textwidth}{@{\extracolsep{\fill}}lrrrrrrrrr@{}}
\toprule
 & \multicolumn{4}{c}{Discrimination $\uparrow$} & \multicolumn{3}{c}{Probability accuracy $\downarrow$} & \multicolumn{2}{c}{Complementary task $\uparrow$} \\
\cmidrule(lr){2-5}\cmidrule(lr){6-8}\cmidrule(lr){9-10}
Design & All & Human & Within & Within & Log loss & Brier & Phase & Winner & C-index \\
 & AUC & AUC & phase & match & & & cal. error & accuracy & \\
\midrule
Ours, GBDT & 0.768 & 0.762 & 0.742 & \textbf{0.747} & 0.498 & \textbf{0.163} & 0.020 & 0.347 & 0.647 \\
Ours, PTR-Former & \textbf{0.769} & \textbf{0.763} & \textbf{0.751} & \textbf{0.747} & \textbf{0.496} & \textbf{0.163} & 0.025 & 0.374 & 0.654 \\
Ours, regularized logistic & 0.759 & 0.752 & 0.737 & 0.729 & 0.504 & 0.166 & 0.019 & 0.366 & 0.646 \\
KRAFTON adaptation & 0.720 & 0.714 & 0.722 & 0.738 & 0.529 & 0.176 & 0.068 & \textbf{0.392} & \textbf{0.657} \\
KRAFTON + phase baseline & 0.747 & 0.740 & 0.722 & 0.738 & 0.513 & 0.169 & \textbf{0.012} & \textbf{0.392} & \textbf{0.657} \\
Sagong adaptation (XGBoost) & 0.675 & 0.669 & 0.620 & 0.627 & 0.575 & 0.193 & 0.018 & --- & --- \\
\midrule
\multicolumn{10}{l}{\emph{KRAFTON backbone, our inputs (reference recipes)}} \\
Weighted survival objective & 0.730 & 0.727 & 0.729 & 0.736 & 0.523 & 0.173 & 0.062 & \textbf{0.372} & \textbf{0.655} \\
Next-phase cross-entropy & \textbf{0.767} & \textbf{0.760} & \textbf{0.747} & \textbf{0.744} & \textbf{0.497} & \textbf{0.163} & \textbf{0.014} & 0.347 & 0.648 \\
\midrule
\multicolumn{10}{l}{\emph{Both input sets combined}} \\
GBDT, ours + KRAFTON inputs & 0.786 & 0.783 & 0.762 & 0.768 & 0.482 & 0.157 & 0.020 & 0.358 & 0.659 \\
\bottomrule
\end{tabular*}
\smallskip
\parbox{\textwidth}{\scriptsize Bold: best rounded value per block; within columns are AUCs. Reference recipes differ in stopping and intercept initialization (matched controls: Table~\ref{tab:objective}). Winner accuracy and C-index use the ranking output (PTR: placement-ranking head).}
\end{table*}

The combined-input GBDT has the highest mean next-phase AUC in the reference matrix (0.786), but it uses additional information, so this is not an equal-information architecture comparison. With our inputs alone, mean AUC is 0.768 (GBDT) and 0.769 (PTR-Former), compared with 0.720 for the KRAFTON adaptation, 0.747 after phase-baseline adjustment, and 0.675 for the Sagong adaptation (Table~\ref{tab:main}). The random-effects pooled PTR advantage over phase-adjusted KRAFTON is $+0.0197$ $[0.0143,\allowbreak 0.0250]$, with substantial heterogeneity ($I^2=0.85$); GBDT's is $+0.0191$ $[0.0138,\allowbreak 0.0244]$. PTR exceeds Sagong by $+0.0922$ $[0.0830,\allowbreak 0.1014]$. The regularized logistic baseline attains mean AUC 0.759; pooled, it trails GBDT by $0.0098$ $[0.0034,\allowbreak 0.0162]$ and exceeds phase-adjusted KRAFTON by $0.0103$ $[0.0009,\allowbreak 0.0196]$, but not within matches ($-0.0088$ $[-0.0203,\allowbreak 0.0028]$).

Fig.~\ref{fig:contrast}(a) compares observed and predicted elimination rates by phase. KRAFTON's phase calibration error decreases from 0.068 to 0.012 after baseline adjustment, versus 0.020 for GBDT and 0.025 for PTR. The adjustment changes probability levels but preserves rankings within phases and matches. Relative to phase-adjusted KRAFTON, within-phase AUC is higher by $+0.0271$ $[0.0180,\allowbreak 0.0361]$ for PTR and $+0.0187$ $[0.0099,\allowbreak 0.0276]$ for GBDT; within-match gains are $+0.0071$ $[0.0030,\allowbreak 0.0113]$ and $+0.0100$ $[0.0071,\allowbreak 0.0128]$. Centering logits within match--phases cuts the within-phase gains to $+0.0093$ $[0.0044,\allowbreak 0.0143]$ and $+0.0073$ $[0.0020,\allowbreak 0.0126]$, close to the within-match gains (Fig.~\ref{fig:contrast}c), whereas match--phase mean logits alone favor PTR and GBDT by $+0.048$ $[0.036,\allowbreak 0.060]$ and $+0.043$ $[0.032,\allowbreak 0.054]$.

The advantage persists among human teams (Table~\ref{tab:main}): the pooled PTR--KRAFTON gap is $+0.0204$ $[0.0141,\allowbreak 0.0266]$, while PTR minus GBDT is essentially zero ($-0.00001$ $[-0.0038,\allowbreak 0.0038]$). AI--human separation alone therefore cannot explain the advantage. Nor can match type: pooled AUC gains over phase-adjusted KRAFTON are $+0.0166$ (normal) and $+0.0246$ (ranked) for PTR and $+0.0157$ and $+0.0315$ for GBDT, in ten and six cohorts with at least 20 test matches of that type (all intervals exclude zero); PTR minus GBDT is uncertain in both.

The survival recipe better predicts winners: KRAFTON exceeds PTR by $0.0175$ $[0.0020,\allowbreak 0.0324]$, while their C-index difference is uncertain ($0.0024$ $[-0.0005,\allowbreak 0.0053]$). With our inputs/backbone fixed, weighted survival exceeds BCE by $0.0252$ winner accuracy and $0.0065$ C-index (both intervals exclude zero).

For retrospective review, the three highest-risk teams at a boundary include 51.7\% of next-phase eliminations under GBDT, 51.6\% under PTR-Former, 51.1\% under phase-adjusted KRAFTON, and 53.3\% with combined inputs, versus 34.7\% at random and 45.1\% for the fewest-remaining-members rule (pooled GBDT gain over this rule $+0.068$ $[0.062,\allowbreak 0.073]$, 11/11). GBDT, PTR-Former, and phase-adjusted KRAFTON lie within 0.6\,pp of one another, about a tenth of their 6.0--6.6\,pp gain over the rule. Mean calibration slopes are 0.86 (GBDT), 0.95 (PTR-Former), and 0.95 (phase-adjusted KRAFTON), where values below one indicate overly spread predictions; GBDT's calibration-in-the-large of $-0.09$ indicates mild overprediction on patch 41.2.

\subsection{RQ3: Decomposing Design Factors}
\label{sec:results:rq3}

\begin{table}[t]
\centering\footnotesize
\caption{Conditional reference-path changes in AUC, pooled across 11 cohorts.}
\label{tab:steps}
\setlength{\tabcolsep}{3pt}
\begin{tabular}{lrrr}
\toprule
Step & Overall & Within phase & Within match \\
\midrule
(a) Sagong inputs $\to$ ours & \textbf{+0.075} & \textbf{+0.106} & \textbf{+0.101} \\
(b) KRAFTON inputs $\to$ ours & -0.003 & \textbf{-0.007} & \textbf{-0.012} \\
(c1) + phase baseline & \textbf{+0.020} & $0^{\dagger}$ & $0^{\dagger}$ \\
(c2) $\to$ next-phase BCE recipe & \textbf{+0.016} & \textbf{+0.021} & \textbf{+0.009} \\
(d) + explicit phase intercepts & -0.0003 & +0.0001 & +0.0004 \\
(e) $\to$ team-axis PTR & +0.004 & +0.005 & +0.004 \\
(f) + phase axis, relation bias & -0.002 & -0.002 & -0.002 \\
(g) + KRAFTON inputs (union) & \textbf{+0.017} & \textbf{+0.019} & \textbf{+0.020} \\
\bottomrule
\end{tabular}
\smallskip
\parbox{\columnwidth}{\scriptsize Random-effects pooled estimates. Bold: 95\% CI excludes 0 and $|\Delta|\geq0.005$. $\dagger$Zero by monotonicity within a phase. Steps (a), (b), (g) use GBDT. Other steps change reference recipes and may bundle stopping, initialization, or baseline-fit split; matched controls follow.}
\end{table}

Table~\ref{tab:steps} pools the path changes with random effects. On the Sagong path, adding habitat context to Sagong's inputs raises GBDT AUC by $+0.046$ $[0.041,\allowbreak 0.052]$, and replacing the rest with our inputs by $+0.028$ $[0.024,\allowbreak 0.033]$ (both 11/11). Removing the habitat block from our GBDT costs $0.0127$ $[0.0110,\allowbreak 0.0144]$ overall AUC and $0.0151$ $[0.0119,\allowbreak 0.0182]$ within-phase AUC (11/11 for both).

On the KRAFTON path, the phase baseline ($+0.020$) and the BCE recipe ($+0.016$) contribute most; these recipe changes bundle horizon, optimization, and mapping choices, so matched controls follow.

\paragraph{Horizon alignment explains the within-match advantage.}
\begin{table}[t]
\centering\footnotesize
\caption{Matched objective and horizon controls. AUCs average three seeds after a common phase-aware validation mapping.}
\label{tab:objective}
\setlength{\tabcolsep}{3pt}
\begin{tabular}{lrrrr}
\toprule
 & \multicolumn{2}{c}{Erangel} & \multicolumn{2}{c}{Taego} \\
\cmidrule(lr){2-3}\cmidrule(lr){4-5}
Loss and horizon & Within & Within & Within & Within \\
 & phase & match & phase & match \\
\midrule
Weighted, full & 0.766 & 0.769 & 0.767 & 0.776 \\
Weighted, next phase & 0.772 & 0.778 & 0.773 & 0.789 \\
Cox, full & 0.764 & 0.766 & 0.762 & 0.772 \\
Cox, next phase & 0.773 & 0.778 & 0.775 & 0.787 \\
BCE, next phase & 0.779 & 0.779 & 0.782 & 0.786 \\
\bottomrule
\end{tabular}
\end{table}

With architecture, inputs, optimization, and stopping aligned, truncating standard Cox at the next-phase boundary improves within-match AUC by $+0.0123$ $[0.0088,\allowbreak 0.0160]$ in Erangel and $+0.0147$ $[0.0112,\allowbreak 0.0180]$ in Taego (Table~\ref{tab:objective}, Fig.~\ref{fig:contrast}b). BCE minus truncated Cox is then only $+0.0008$ $[-0.0013,\allowbreak 0.0028]$ and $-0.0008$ $[-0.0028,\allowbreak 0.0012]$. Thus the within-match advantage over full-follow-up objectives largely reflects horizon alignment. BCE retains a within-phase advantage over truncated Cox, $+0.0060$ $[0.0038,\allowbreak 0.0082]$ and $+0.0076$ $[0.0055,\allowbreak 0.0097]$, and the diagnostics place it at the match level: same-match pairs are only 0.2\% of within-phase pairs, centering logits within each match--phase shrinks the gap to $+0.0015$ $[-0.0003,\allowbreak 0.0033]$ and $+0.0008$ $[-0.0009,\allowbreak 0.0024]$, and match--phase mean logits alone rank matches better under BCE by $+0.016$ $[0.013,\allowbreak 0.020]$ and $+0.019$ $[0.015,\allowbreak 0.022]$, although truncated Cox's means still reach AUC 0.586 and 0.597 because team inputs carry match state, about 84\% of BCE's above-chance match-level discrimination (Fig.~\ref{fig:contrast}c). Truncating Cox leaves this match-level component unchanged ($+0.002$ and $-0.004$).

\paragraph{A two-factor attribution, averaged over orders.}
The objective--mapping factorial yields an equal-cohort mean AUC gain of $0.0358$ $[0.0336,\allowbreak 0.0380]$. The mapping's Shapley contribution is $0.0108$ $[0.0098,\allowbreak 0.0118]$ and the objective's is $0.0250$ $[0.0234,\allowbreak 0.0265]$. The mapping contributes zero within-phase AUC, whereas the objective contributes $0.0137$ $[0.0117,\allowbreak 0.0157]$. The objective factor includes the horizon.

\paragraph{Structural effects depend on metric and threshold.}
\begin{table}[t]
\centering\footnotesize
\caption{Five primary structure decisions: random-effects estimates over 11 cohorts.}
\label{tab:primary}
\setlength{\tabcolsep}{2.5pt}
\begin{tabular}{lrrl}
\toprule
Comparison & Estimate & 90\% CI & Decision \\
\midrule
Phase axis, AUC & -0.0013 & [-0.0024, -0.0001] & Equiv. \\
Team axis, AUC & +0.0043 & [+0.0031, +0.0055] & Inconcl. \\
Team axis, rank $\rho$ & +0.0083 & [+0.0028, +0.0137] & Inconcl. \\
Relation alignment, AUC & +0.0004 & [+0.0002, +0.0005] & Equiv. \\
Zone token, AUC & +0.0015 & [+0.0006, +0.0024] & Equiv. \\
\bottomrule
\end{tabular}
\smallskip
\parbox{\columnwidth}{\scriptsize Bounds are $\pm0.005$ for AUC and $\pm0.01$ for rank correlation. Holm-adjusted one-sided $p$ against zero, in row order: 0.965, $2.22\times10^{-8}$, 0.0122, 0.00259, 0.00924. Equivalence uses unadjusted 90\% intervals.}
\end{table}

Three primary comparisons meet practical equivalence; both team-axis comparisons are inconclusive (Table~\ref{tab:primary}). REML/Hartung--Knapp retains these verdicts. Team-axis AUC is positive (Holm $p<0.001$), but its 90\% interval crosses 0.005; the rank-correlation interval crosses 0.01. The team-axis AUC verdict would be a gain for $\delta\le0.0043$ and equivalence for $\delta\ge0.0055$; we retain $\delta=0.005$. Across nine design contrasts in 11 cohorts, top-3 capture changes by 0.66 $[0.54,\allowbreak 0.78]$ per unit of within-match AUC, so a within-match AUC difference of 0.005 corresponds to about 0.33\,pp, or one more captured elimination per about 90 boundaries; the primary structure contrasts change capture by at most 0.36\,pp.

Team attention improves within-phase AUC by $+0.0108$ $[0.0086,\allowbreak 0.0130]$ (11/11). PTR exceeds GBDT by $+0.0058$ $[0.0035,\allowbreak 0.0080]$ on this metric, but their within-match difference is uncertain ($-0.001$ $[-0.006,\allowbreak 0.004]$). The centering diagnostic places these gains across matches: after centering, PTR's within-phase advantage over GBDT falls to $+0.0011$ $[-0.0022,\allowbreak 0.0044]$, and match--phase mean logits alone favor PTR by $+0.0034$ $[0.0016,\allowbreak 0.0051]$; for the team axis, $+0.0049$ $[0.0022,\allowbreak 0.0076]$ is match-level and $+0.0043$ $[0.0020,\allowbreak 0.0067]$ remains after centering. Relation alignment contributes only $+0.00035$ overall AUC ($I^2=0.93$; Hartung--Knapp interval crossing zero).

\paragraph{The input sets are complementary.}
With GBDT fixed, our inputs score slightly below the reconstructed KRAFTON inputs: pooled AUC difference $-0.0035$ $[-0.0066,\allowbreak -0.0004]$, with practical equivalence unresolved. The deficit is larger within matches ($-0.0124$ $[-0.0170,\allowbreak -0.0078]$). Combining the input sets improves AUC over ours by $+0.0170$ $[0.0141,\allowbreak 0.0198]$ (11/11).

\paragraph{Equal-trial tuning narrows architecture differences.}
With common inputs and BCE, the overall AUC difference between PTR and the KRAFTON backbone is $-0.0011$ in Erangel and $+0.0006$ in Taego; both intervals include zero (Table~\ref{tab:tuning}). The within-match differences from GBDT are $-0.0062$ $[-0.0097,\allowbreak -0.0029]$ and $+0.0026$ $[-0.0009,\allowbreak 0.0059]$. Against the reference recipes, tuning on patch 41.1 alone with 153 columns (adding the centroid coordinates) changes GBDT and PTR AUC by at most 0.002 (Table~\ref{tab:tuning} note). The KRAFTON backbone gains 0.003/0.007, which removes Taego's reference gap to PTR ($+0.006$ to $+0.0006$); step (e)'s $+0.004$ in Table~\ref{tab:steps} also lies within this recipe range. No architecture dominates both cohorts and metrics under this search budget.
\begin{table}[t]
\centering\footnotesize
\caption{Matched-input, equal-trial tuning. Each value averages three selected-model seeds; all use the BCE target.}
\label{tab:tuning}
\setlength{\tabcolsep}{3pt}
\begin{tabular}{lrrrr}
\toprule
Model & All AUC & Human & Within & Log loss \\
 & & AUC & match & \\
\midrule
\multicolumn{5}{l}{\emph{Erangel TPP squad}} \\
GBDT & 0.803 & 0.804 & 0.783 & 0.452 \\
KRAFTON-derived & 0.803 & 0.803 & 0.779 & 0.453 \\
PTR (no rank loss) & 0.801 & 0.802 & 0.777 & 0.454 \\
\multicolumn{5}{l}{\emph{Taego TPP squad}} \\
GBDT & 0.805 & 0.805 & 0.789 & 0.443 \\
KRAFTON-derived & 0.808 & 0.808 & 0.789 & 0.442 \\
PTR (no rank loss) & 0.809 & 0.808 & 0.791 & 0.442 \\
\bottomrule
\end{tabular}
\smallskip
\parbox{\columnwidth}{\scriptsize Reference recipes (Table~\ref{tab:main}; 151 inputs, PTR-Former with ranking loss) reach all-team AUC 0.803/0.803 (GBDT), 0.799/0.807 (PTR-Former), and 0.800/0.801 (KRAFTON backbone, cross-entropy) in Erangel/Taego.}
\end{table}


\paragraph{Population, history, and graph controls.}
Population and history controls move AUC by at most 0.0033: removing AI input changes human AUC by $-0.0012$ to $+0.0010$ and human-only loss by $-0.0022$ to $+0.0018$ across six model--cohort pairs, past-token summaries lower KRAFTON AUC by 0.0033/0.0028, and ordered PTR prefixes beat shuffled ones by at most 0.0010. GATv2 (AUC 0.8048/0.8079 in Erangel/Taego) exceeds PTR by $+0.0034$/$+0.0023$ but not the KRAFTON backbone ($+0.0009$/$+0.0003$, uncertain), relation alignment adds $+0.0020$/$+0.0010$ over jointly relabeled graphs, and every graph effect stays below the 0.005 threshold.



%======================================================================
\section{Discussion}
\label{sec:discussion}
\subsection{Align the Target Before Attributing Gains to Structure}
\label{sec:discussion:drivers}
Next-phase truncation leaves no clear within-match difference between Cox and BCE, so horizon alignment is a major contributor. The remaining within-phase advantage is consistent with a structural explanation. Cox-type objectives normalize scores within each match's risk set: adding a constant to every score in a match--phase leaves the loss unchanged, so they give no training signal for a match-level risk offset, and match information enters Cox scores only as a by-product of team-level fitting. BCE instead responds to such a shift: its derivative with respect to a common offset $c_{m,k}$ is $\sum_{i\in\mathcal{R}_{m,k}}(\hat p_{m,i,k}-Y_{m,i,k})$, the gap between expected and observed eliminations, so it provides a training signal for match-level risk wherever the model can represent such offsets. Because within-phase AUC mostly compares teams from different matches, the difference appears within phases but not within matches, and centering substantially attenuates it, leaving no clear residual difference in the matched Cox--BCE comparison. Conditional likelihoods in matched and fixed-effects designs share this property: group-level effects drop out~\cite{Chamberlain80}. Learning to rank faces the same invariance: without ties, the Cox partial likelihood is the Plackett--Luce likelihood that ListMLE maximizes~\cite{Xia08}, such listwise scores are uncalibrated across queries, and adding a pointwise sigmoid cross-entropy is one remedy~\cite{Bai23}. What we add is a matched-control measurement of this component in landmark survival prediction. Phase-specific mapping separately improves probability accuracy while preserving within-phase rankings.

Removing habitat inputs worsens prediction. KRAFTON inputs outperform ours within matches under GBDT; their union improves further, indicating complementary information. A linear model on our inputs recovers about half of GBDT's overall gain over phase-adjusted KRAFTON but none of its within-match gain, so the within-match advantage of next-phase learning appears only with a nonlinear learner.

Team attention and PTR-Former raise within-phase but not within-match AUC, and the same centering diagnostic places PTR-Former's gain over GBDT, and about half of the team-axis gain, in match-level context. GATv2 detects small alignment effects and scores close to KRAFTON; existing relation summaries may limit the value of explicit graphs.

\subsection{Interpreting Habitat Associations}
Competition is strongest early and inside depth and members later; the depth pattern resembles an ecological edge effect~\cite{Ries04} without identifying its mechanism. The population-sensitive movement residual and the lite model's 73\% retained within-phase discrimination support cautious, non-causal interpretation.

\subsection{Uses and Limitations}
\label{sec:discussion:limits}
The supported application is retrospective spectator-view match review; top-3 capture is only a proxy, analyst benefit and latency are unmeasured, and live use would require an online eligibility audit and real-time feeds. The learned designs beat the fewest-remaining-members rule by about 6\,pp in top-3 capture but differ by at most 0.6\,pp among themselves, so the study informs evaluation practice more than the choice of a review system.

Eligibility uses final elimination times, which future returns can affect, so online risk-set validity is unverified; the cutoff checks cover retained rows only. The data span ten weeks and three patches with one solo cohort, Karakin has six training matches, and there is no prospective or new-rule test; cohorts share a collection window, and bootstrap intervals condition on fitted models. Baselines use proxy inputs without original checkpoints and change Sagong's unit and interval, so adaptation quality can affect their comparisons; without redistributed telemetry, third parties can inspect the code but not rerun the exact samples. Habitat associations are observational: teams choose positions, and zone and combat damage are not separated. Tuning covers two cohorts, a limited search, and three seeds without equalizing computation; a standalone MLP baseline remains untested, and ranked-match strata cover six cohorts with 26--108 test matches each.

\paragraph{Code availability.}
Anonymized training, evaluation, and figure-generation code, feature definitions, configurations, and verification notes are available at \artifacturl{}; raw and processed telemetry are not redistributed.

\section{Conclusion}
\label{sec:conclusion}
Patch-separated evaluation identifies input representation, probability mapping, and prediction horizon as major contributors to next-phase elimination prediction. Truncated Cox closes BCE's within-match advantage; BCE's remaining within-phase advantage, like most of the within-phase gain over phase-adjusted KRAFTON, is consistent with match-level risk offsets, for which risk-set-normalized objectives provide no training signal. Combining the two input sets gives the highest mean next-phase AUC in the reference matrix; equal-trial comparisons in two cohorts yield no consistent architecture winner under the tested search budget. Competition and members show consistent observational associations; movement residuals are population-sensitive. Future work should test combined Cox and cross-entropy objectives, private-observation inputs, and newly collected patches, including the zone-rule revamp~\cite{PUBGBlueZone2026}, with additional assumptions required for causal interpretation.

\section*{AI Assistance Disclosure}
Anthropic's Claude and OpenAI Codex assisted with writing and revising the analysis and figure-generation code; the authors reviewed all code, numbers, and text and take responsibility for the content.

\bibliographystyle{ACM-Reference-Format-citation-order}
\bibliography{refs}
\end{document}
